%% lalonde.tex — Zurich Day 1 walkthrough: LaLonde (1986), Dehejia-Wahba (2002), the homework
%% Scott Cunningham, University of Zurich, May 12, 2026

\documentclass[11pt,aspectratio=169]{beamer}
\usepackage{remix}
\usepackage{booktabs}
\renewcommand{\remixfooter}{University of Zurich\,\textperiodcentered\,Department of Economics\,\textperiodcentered\,May 2026}

\begin{document}

%% ---------- Title ----------
\remixtitleframe
  {LaLonde, Dehejia--Wahba, and Why the 2$\times$2 Still Matters}
  {The story behind yesterday's homework}
  {Scott Cunningham \textperiodcentered\ Baylor University}
  {University of Zurich \textperiodcentered\ May 12, 2026}

%% =====================================================================
%%                                LESSON 1
%% =====================================================================
\remixsection{1.\ LaLonde 1986}

\begin{frame}{The National Supported Work Demonstration}
\vspace{1em}
\begin{center}
\large
A federally funded \textbf{randomized} job-training program.
\end{center}
\vspace{0.4em}
\begin{center}
\large
1975--1979. Four target groups, including \\[0.2em] disadvantaged adult men.
\end{center}
\vspace{0.6em}
\begin{center}
\small\color{warmgray}\itshape
Random assignment to treatment vs.\ control. \;Outcome: real earnings in 1978.
\end{center}
\end{frame}

\begin{frame}{LaLonde's experimental benchmark}
\vspace{1.6em}
\begin{center}
\Large
$\widehat{\text{ATT}}^{\text{RCT}} \;\approx\; \$886$
\end{center}
\vspace{0.6em}
\begin{center}
1978 dollars. \;Adult-male subsample.
\end{center}
\vspace{1em}
\begin{center}
\small\color{warmgray}\itshape
Random assignment $\Rightarrow$ no selection. \;This is the \textit{truth}.
\end{center}
\end{frame}

\begin{frame}{Then LaLonde threw away the controls}
\vspace{1em}
\begin{center}
\large
He kept the treated --- the 297 NSW men.
\end{center}
\vspace{0.6em}
\begin{center}
\large
He replaced the controls with samples from the \\[0.2em] \textbf{CPS} and the \textbf{PSID}.
\end{center}
\vspace{0.8em}
\begin{center}
\small\color{warmgray}\itshape
A simulation of what an applied researcher would face --- \\
no experiment, just observational data and methods.
\end{center}
\end{frame}

\begin{frame}{Then he tried every method going at the time}
\vspace{0.6em}
\begin{itemize}
  \item OLS with controls
  \item Two-step Heckman selection corrections
  \item Fixed effects on pre-program earnings
  \item Difference-in-differences
\end{itemize}
\vspace{0.6em}
\begin{center}
\large\itshape\color{charcoal}
Could any of them recover \$886?
\end{center}
\end{frame}

\begin{frame}{No}
\vspace{1.6em}
\begin{center}
\large
Estimates ranged from \;\highlight{$-$\$16{,}000} \;to \;\highlight{$+$\$3{,}000}.
\end{center}
\vspace{1em}
\begin{center}
\large
Most were negative. \;Most missed the benchmark by orders of magnitude.
\end{center}
\vspace{1em}
\begin{center}
\small\color{warmgray}\itshape
LaLonde (1986), \textit{AER} 76(4): 604--620.
\end{center}
\end{frame}

\begin{frame}{The headline that shook applied econometrics}
\vspace{1.6em}
\begin{center}
\Large\itshape\color{charcoal}
``The econometric procedures \\[0.2em] cannot replicate the experimental results.''
\end{center}
\vspace{1.4em}
\begin{center}
\small\color{warmgray}\itshape
This is the paper that made program evaluation \\
move toward randomization and design.
\end{center}
\end{frame}

%% =====================================================================
%%                                LESSON 2
%% =====================================================================
\remixsection{2.\ Dehejia \& Wahba 2002}

\begin{frame}{Dehejia \& Wahba reopened the case}
\vspace{1em}
\begin{center}
\large
\textbf{Rajeev Dehejia} and \textbf{Sadek Wahba}, \\[0.2em] \textit{Review of Economics and Statistics} (2002).
\end{center}
\vspace{0.8em}
\begin{center}
\large
They asked: \;does \textbf{propensity score matching} \\[0.2em] do better than LaLonde's methods?
\end{center}
\end{frame}

\begin{frame}{First, they re-cut the experimental sample}
\vspace{1em}
\begin{center}
\large
They restricted to NSW men with earnings observed in \\[0.2em] \textbf{both 1974 and 1975}.
\end{center}
\vspace{0.6em}
\begin{center}
\large
A smaller, slightly different population.
\end{center}
\vspace{1em}
\begin{center}
\small\color{warmgray}\itshape
This is the subsample the homework uses: \texttt{lalonde\_exp\_panel.dta}.
\end{center}
\end{frame}

\begin{frame}{A new experimental benchmark}
\vspace{1.6em}
\begin{center}
\Large
$\widehat{\text{ATT}}^{\text{RCT}}_{\text{D\&W}} \;\approx\; \$1{,}794$
\end{center}
\vspace{0.6em}
\begin{center}
\large
Roughly \;\highlight{\$1{,}700} \;in the way we usually report it.
\end{center}
\vspace{1em}
\begin{center}
\small\color{warmgray}\itshape
Why bigger than LaLonde's \$886? \;\textit{Different subsample, different ATT.}
\end{center}
\end{frame}

\begin{frame}{Why the headlines differ}
\vspace{0.5em}
\begin{itemize}
  \item Both numbers come from the \textit{same} NSW experiment.
  \item LaLonde used the full adult-male sample.
  \item D\&W used a \textit{subset} --- only men with two years of pre-program earnings.
  \item Subsetting on a covariate changes the population whose ATT you're estimating.
  \item Same identification (random assignment). \;Different target group.
\end{itemize}
\vspace{0.5em}
\begin{center}
\small\color{warmgray}\itshape
This is not a contradiction --- it is the ATT defined on different units.
\end{center}
\end{frame}

\begin{frame}{D\&W's quasi-experimental result}
\vspace{1em}
\begin{center}
\large
With \textbf{propensity-score matching} on \\[0.2em] CPS / PSID controls...
\end{center}
\vspace{0.6em}
\begin{center}
\Large
they recovered \;\highlight{$\approx \$1{,}700$}.
\end{center}
\vspace{0.8em}
\begin{center}
\small\color{warmgray}\itshape
A constructive answer to LaLonde: \;design plus matching \textit{can} work.
\end{center}
\end{frame}

%% =====================================================================
%%                                LESSON 3
%% =====================================================================
\remixsection{3.\ What you did yesterday}

\begin{frame}{Yesterday's lab was on D\&W's experimental subsample}
\vspace{1em}
\begin{center}
\large
\texttt{lalonde\_exp\_panel.dta} \;--- a panel of NSW men, \\[0.2em] years 74, 75, 78.
\end{center}
\vspace{0.6em}
\begin{center}
\large
\textbf{The benchmark (post-period diff): $\$1{,}794$.}
\end{center}
\vspace{0.6em}
\begin{center}
\small\color{warmgray}\itshape
You had \textit{no covariates}. \;Just \texttt{re}, \texttt{ever\_treated}, \texttt{year}. \;\;$n_T = 185, \;n_C = 260$.
\end{center}
\end{frame}

\begin{frame}{Step 1: four averages, by hand}
\vspace{0.4em}
\begin{center}
\small
\renewcommand{\arraystretch}{1.4}
\begin{tabular}{@{}lcc@{}}
\toprule
 & \textbf{Pre (75)} & \textbf{Post (78)} \\
\midrule
Treated   & $\bar Y^T_{75}$ & $\bar Y^T_{78}$ \\
Control   & $\bar Y^C_{75}$ & $\bar Y^C_{78}$ \\
\bottomrule
\end{tabular}
\end{center}
\vspace{0.8em}
\begin{center}
\small\color{warmgray}\itshape
Four averages. Three subtractions. \;One DiD.
\end{center}
\end{frame}

\begin{frame}{Step 2: the DiD}
\vspace{1em}
\begin{center}
$\widehat{\delta}_{\text{DiD}} \;=\; (\,\bar Y^T_{78} - \bar Y^T_{75}\,) \;-\; (\,\bar Y^C_{78} - \bar Y^C_{75}\,)$
\end{center}
\vspace{0.6em}
\begin{center}
\small
$=\; (\$6{,}349 - \$1{,}532) \;-\; (\$4{,}555 - \$1{,}267) \;=\; \$4{,}817 - \$3{,}288$
\end{center}
\vspace{0.6em}
\begin{center}
\Large
$\widehat{\delta}_{\text{DiD}} \;=\; \$1{,}529$
\end{center}
\vspace{0.6em}
\begin{center}
\small\color{warmgray}\itshape
Post-period diff-in-means alone: \;$\$6{,}349 - \$4{,}555 = \$1{,}794$.
\end{center}
\end{frame}

\begin{frame}{Two valid estimators, two different numbers}
\vspace{0.8em}
\begin{center}
\small
\renewcommand{\arraystretch}{1.3}
\begin{tabular}{@{}lr@{}}
\toprule
Post-period diff-in-means \;(78 only)        & \$1{,}794 \\
DiD \;(78 vs.\ 75)                            & \$1{,}529 \\
\midrule
Difference                                    & \$\phantom{1{,}}265 \\
\bottomrule
\end{tabular}
\end{center}
\vspace{0.6em}
\begin{center}
\large
The pre-period gap was \;$\bar Y^T_{75} - \bar Y^C_{75} = +\$265$.
\end{center}
\vspace{0.6em}
\begin{center}
\large
DiD subtracts it. \;Post-period diff doesn't.
\end{center}
\end{frame}

\begin{frame}{Unbiasedness is a repeated-sampling concept}
\vspace{1em}
\begin{center}
\large
Random assignment makes \;$\mathbb{E}[\widehat\delta] \;=\; \text{ATT}$
\end{center}
\vspace{0.6em}
\begin{center}
\large
across hypothetical re-randomizations of the experiment.
\end{center}
\vspace{1em}
\begin{center}
\large
It does \textit{not} make \;$\widehat\delta \;=\; \text{ATT}$ \;in any one sample.
\end{center}
\vspace{0.6em}
\begin{center}
\small\color{warmgray}\itshape
Our pre-period happened to have $+\$265$ of random imbalance. \;Different draw, different imbalance, different DiD.
\end{center}
\end{frame}

\begin{frame}{Step 3: three regressions, same number}
\vspace{0.6em}
\begin{itemize}
  \item \textbf{OLS with interactions}: \;\texttt{re \textasciitilde\ ever\_treated * post}
  \item \textbf{TWFE}: \;\texttt{feols(re \textasciitilde\ ever\_treated:post | id + year)}
  \item \textbf{Long differences}: \;\texttt{lm(re78 $-$ re75 \textasciitilde\ ever\_treated)}
\end{itemize}
\vspace{0.6em}
\begin{center}
\large
All three return \;\highlight{$\widehat\delta = \$1{,}529$}, to machine precision.
\end{center}
\vspace{0.6em}
\begin{center}
\small\color{warmgray}\itshape
The equivalence breaks once we add staggered timing, covariates, or weights. \;Tomorrow.
\end{center}
\end{frame}

\begin{frame}{Step 4: the event study tells the same story}
\vspace{0.4em}
\begin{itemize}
  \item $\beta_{74} = -\$277$ \;--- small, close to zero. \;\highlight{Pre-trend check passes.}
  \item $\beta_{75}$ is the omitted baseline (zero by construction).
  \item $\beta_{78} = +\$1{,}529$ \;--- same as the DiD.
\end{itemize}
\vspace{0.6em}
\begin{center}
\small\color{warmgray}\itshape
Each event-study coefficient is a 2$\times$2 DiD with 75 as baseline. \;Same arithmetic, three points.
\end{center}
\end{frame}

%% =====================================================================
%%                                LESSON 4
%% =====================================================================
\remixsection{4.\ Why the ATT didn't move}

\begin{frame}{A puzzle}
\vspace{1em}
\begin{center}
\large
We used \textbf{non-experimental controls} (CPS, PSID) \\[0.2em] in the larger LaLonde--D\&W literature.
\end{center}
\vspace{0.6em}
\begin{center}
\large
And yet the ATT we target is still \;\highlight{$\approx \$1{,}700$.}
\end{center}
\vspace{0.8em}
\begin{center}
\large\itshape\color{charcoal}
Why doesn't the ATT change when we change the controls?
\end{center}
\end{frame}

\begin{frame}{Because the ATT is a property of the treated}
\vspace{1.4em}
\begin{center}
\Large
$\text{ATT} \;=\; \mathbb{E}[Y^1 - Y^0 \mid D = 1,\, \text{Post}]$
\end{center}
\vspace{1.2em}
\begin{center}
\large
Only the \;\highlight{$D = 1$} \;side appears.
\end{center}
\vspace{0.6em}
\begin{center}
The control units are just a \textit{tool} to learn about $\mathbb{E}[Y^0 \mid D{=}1, \text{Post}]$.
\end{center}
\end{frame}

\begin{frame}{Random assignment gave us a clean tool}
\vspace{1.4em}
\begin{center}
\large
Under random assignment:
\end{center}
\vspace{0.4em}
\begin{center}
\Large
$\mathbb{E}[Y^0 \mid D{=}1] \;=\; \mathbb{E}[Y^0 \mid D{=}0]$
\end{center}
\vspace{1em}
\begin{center}
So $\mathbb{E}[Y \mid D{=}0, \text{Post}]$ \textit{is} the counterfactual we need. \\[0.2em]
No selection bias.
\end{center}
\end{frame}

\begin{frame}{Non-experimental controls break that equality}
\vspace{1em}
\begin{center}
\large
CPS men and NSW participants are not interchangeable.
\end{center}
\vspace{0.6em}
\begin{center}
\large
$\mathbb{E}[Y^0 \mid D{=}1] \;\neq\; \mathbb{E}[Y^0 \mid D{=}0]$
\end{center}
\vspace{1em}
\begin{center}
\large
Selection bias enters.
\end{center}
\vspace{0.6em}
\begin{center}
\small\color{warmgray}\itshape
NSW men had worse pre-program earnings. \;The CPS pool didn't.
\end{center}
\end{frame}

\begin{frame}{But the \textit{target} hasn't changed}
\vspace{1.2em}
\begin{center}
\large
$\mathbb{E}[Y^0 \mid D{=}1, \text{Post}]$ is a feature of the NSW men.
\end{center}
\vspace{0.6em}
\begin{center}
\large
It does not depend on who we picked as controls.
\end{center}
\vspace{1em}
\begin{center}
\large
What the controls control is \textit{our estimate} of that target --- \\[0.2em] and the size of the bias.
\end{center}
\end{frame}

\begin{frame}{So what is DiD for, in this story?}
\vspace{0.6em}
\begin{itemize}
  \item In the RCT sample: DiD just \textit{confirms} the post-period difference-in-means.
  \item In the non-experimental sample: DiD tries to \textit{net out} the pre-existing level gap between NSW and CPS men.
  \item The truth is fixed at \$1{,}700. \;The question is whether DiD's parallel-trends assumption is enough to recover it.
\end{itemize}
\vspace{0.6em}
\begin{center}
\large\itshape\color{charcoal}
Did it?
\end{center}
\end{frame}

\begin{frame}{In yesterday's experimental lab: yes (in expectation)}
\vspace{1em}
\begin{center}
\Large
$\widehat{\delta}_{\text{DiD}} \;=\; \$1{,}529$ \;\;vs.\;\; benchmark \;$\$1{,}794$
\end{center}
\vspace{0.8em}
\begin{center}
\large
Within sampling variation. \;Pre-trend was flat. \\[0.3em] Design did its job.
\end{center}
\vspace{1em}
\begin{center}
\small\color{warmgray}\itshape
This afternoon, with covariates and non-experimental controls --- a harder test.
\end{center}
\end{frame}

%% =====================================================================
%%                                LESSON 5
%% =====================================================================
\remixsection{5.\ Now make it harder}

\begin{frame}{Throw away the experimental controls}
\vspace{1em}
\begin{center}
\large
Keep the 297 NSW treated men.
\end{center}
\vspace{0.6em}
\begin{center}
\large
Replace the controls with a random sample of \\[0.2em] American households from the \textbf{CPS}.
\end{center}
\vspace{1em}
\begin{center}
\small\color{warmgray}\itshape
Data: \texttt{lalonde\_nonexp\_panel.dta}. \;Same treated rows. \;New controls.
\end{center}
\end{frame}

\begin{frame}{Selection is severe}
\vspace{0.8em}
\begin{center}
\small
\renewcommand{\arraystretch}{1.3}
\begin{tabular}{@{}lrrr@{}}
\toprule
        & \textbf{1974}    & \textbf{1975}    & \textbf{1978}    \\
\midrule
NSW (treated, $n=185$)   & \$2{,}096        & \$1{,}532        & \$6{,}349        \\
CPS (control, $n=15{,}992$) & \$14{,}017       & \$13{,}651       & \$14{,}847       \\
\midrule
Diff                      & $-\$11{,}921$    & $-\$12{,}119$    & $-\$8{,}498$    \\
\bottomrule
\end{tabular}
\end{center}
\vspace{0.5em}
\begin{center}
\large
$\mathbb{E}[Y^0 \mid D{=}1] \;\ll\; \mathbb{E}[Y^0 \mid D{=}0]$
\end{center}
\vspace{0.3em}
\begin{center}
\small\color{warmgray}\itshape
The level gap is huge --- and so is the \textit{growth-rate} gap. \;DiD doesn't fix that.
\end{center}
\end{frame}

\begin{frame}{Naive DiD on this sample}
\vspace{1.2em}
\begin{center}
\Large
$\widehat{\delta}_{\text{naive DiD}} \;=\; +\$3{,}621$
\end{center}
\vspace{0.6em}
\begin{center}
\large
\textit{Right sign. \;Wrong magnitude.} \;The benchmark is \;\highlight{$\$1{,}794$}.
\end{center}
\vspace{0.8em}
\begin{center}
\small\color{warmgray}\itshape
NSW men had unusually low pre-earnings ($\$1{,}532$) and bounced back fast. \\
CPS controls didn't. \;DiD attributes the gap in growth rates entirely to treatment.
\end{center}
\end{frame}

%% =====================================================================
%%                                LESSON 6
%% =====================================================================
\remixsection{6.\ Four estimators to rescue the ATT}

\begin{frame}{The roster}
\vspace{0.5em}
\begin{itemize}
  \item \textbf{DiD with additive $X$} (your published spec) --- one slope per covariate
  \item \textbf{Outcome Regression} (Heckman, Ichimura \& Todd 1997) --- model the trend on controls; impute the treated counterfactual
  \item \textbf{Inverse Propensity Weighting} (Abadie 2005) --- reweight controls to look like the treated
  \item \textbf{Doubly Robust DiD} (Sant'Anna \& Zhao 2020) --- OR + IPW combined
\end{itemize}
\vspace{0.4em}
\begin{center}
\small\color{warmgray}\itshape
Each lives or dies on a different assumption. \;Compare to the \$1{,}794 benchmark.
\end{center}
\end{frame}

\begin{frame}[fragile]{DiD with additive $X$}
\vspace{0.6em}
\begin{center}
\small
\verb|reg re ever_treated##post X1 X2 ... Xk, robust|
\end{center}
\vspace{0.6em}
\begin{itemize}\small
  \item Pooled OLS on all years. \;No fixed effects.
  \item Covariates enter \textit{linearly} with one slope --- same pre and post.
  \item Identifies the ATT only if Assumptions 1--6 hold (parallel trends, no anticipation, SUTVA, homogeneous $\tau$ in $X$, parallel $X$-trends for treated and control).
\end{itemize}
\vspace{0.4em}
\begin{center}
\small\color{warmgray}\itshape
This is what \texttt{lalonde-did-stata.do} actually computes on the non-experimental panel.
\end{center}
\end{frame}

\begin{frame}{Outcome Regression --- impute the counterfactual}
\vspace{0.6em}
Use the \textit{controls only} to model the trend:
\vspace{0.2em}
\begin{center}
\small
$\widehat\mu_0(X) \;=\; \widehat{\mathbb{E}}[\Delta Y \mid D{=}0,\, X]$
\end{center}
\vspace{0.6em}
Then impute the counterfactual change for each treated unit, and average:
\vspace{0.2em}
\begin{center}
\small
$\widehat\delta_{\text{OR}} \;=\; \dfrac{1}{N_1}\displaystyle\sum_{i:\, D_i=1}\big(\Delta Y_i - \widehat\mu_0(X_i)\big)$
\end{center}
\vspace{0.4em}
\begin{center}
\small\color{warmgray}\itshape
Identification rests on $\widehat\mu_0(X)$ being correctly specified.
\end{center}
\end{frame}

\begin{frame}{Inverse Propensity Weighting}
\vspace{0.8em}
Estimate $\widehat p(X) = \Pr(D{=}1 \mid X)$ by logit. Weight controls by:
\vspace{0.2em}
\begin{center}
$w_i \;=\; \dfrac{\widehat p(X_i)}{1 - \widehat p(X_i)}$
\end{center}
\vspace{0.6em}
\begin{center}
\small
$\widehat\delta_{\text{IPW}} \;=\; \mathbb{E}[\Delta Y \mid D{=}1] \;-\; \mathbb{E}_w[\Delta Y \mid D{=}0]$
\end{center}
\vspace{0.6em}
\begin{center}
\small\color{warmgray}\itshape
Identification rests on $\widehat p(X)$ being correctly specified.
\end{center}
\end{frame}

\begin{frame}{Doubly Robust DiD}
\vspace{1.2em}
\begin{center}
\large
If \textit{either} $\widehat\mu_0(X)$ \textit{or} $\widehat p(X)$ is correct, DR is consistent.
\end{center}
\vspace{1em}
\begin{center}
\large
Use the pscore as a \textbf{weight,} not as a \textbf{regressor.}
\end{center}
\vspace{0.8em}
\begin{center}
\small\color{warmgray}\itshape
Sant'Anna \& Zhao (2020). \;Implemented in \texttt{drdid} (Stata, R).
\end{center}
\end{frame}

%% =====================================================================
%%                                LESSON 7
%% =====================================================================
\remixsection{7.\ How close did each one get?}

\begin{frame}{The comparison}
\vspace{0.3em}
\begin{center}
\footnotesize
\renewcommand{\arraystretch}{1.25}
\begin{tabular}{@{}lrr@{}}
\toprule
\textbf{Estimator (Scott's published \texttt{drdid all})} & \textbf{ATT} & \textbf{Distance from \$1{,}794} \\
\midrule
Naive DiD (no $X$)                                & $+\$3{,}621$ & overshoots by \$1{,}827 \\
DiD with additive $X$ (your \texttt{reg} spec)    & $+\$3{,}522$ & overshoots by \$1{,}728 \\
\midrule
\multicolumn{3}{l}{\textit{Design-aware estimators (\texttt{drdid all}):}} \\
\;\;OR --- Regression-adjusted (\texttt{reg})         & $+\$1{,}770$ & $-\$24$ \\
\;\;IPW --- Abadie (\texttt{ipw})                     & $+\$1{,}861$ & $+\$67$ \\
\;\;DR-IPW (\texttt{dripw})                            & $+\$1{,}979$ & $+\$185$ \\
\;\;DR-improved (\texttt{drimp})                       & $+\$2{,}033$ & $+\$239$ \\
\midrule
\textbf{RCT benchmark (post-period diff)}        & \textbf{$+\$1{,}794$} & --- \\
\bottomrule
\end{tabular}
\end{center}
\vspace{0.3em}
\begin{center}
\footnotesize\color{warmgray}\itshape
Naive DiD and additive-$X$ DiD both overshoot by $\sim\$1{,}800$ --- additive controls barely help. \\[0.2em]
The \texttt{drdid} design-aware estimators all land within $\$240$ of the benchmark; OR is closest at \$24 below.
\end{center}
\end{frame}

\begin{frame}{OR lands closest; all four design-aware estimators are within \$240}
\vspace{0.5em}
\begin{itemize}
  \item \textbf{OR} (\texttt{reg}, \$1{,}770) sits \textbf{\$24 below} --- closest of all.
  \item \textbf{IPW} (Abadie, \$1{,}861) sits \$67 above.
  \item \textbf{DR-IPW} (\$1{,}979) sits \$185 above.
  \item \textbf{DR-improved} (\$2{,}033) sits \$239 above --- furthest of the design-aware row.
  \item Additive-$X$ DiD (\$3{,}522) barely moves off naive DiD (\$3{,}621).
\end{itemize}
\vspace{0.4em}
\begin{center}
\small\color{warmgray}\itshape
``Doubly robust'' means \textit{robust to one misspecification,} not magically best in every sample. \\
With $n_T = 185$ and $n_C \approx 16{,}000$, sampling variation alone explains the spread among design-aware estimators.
\end{center}
\end{frame}

\begin{frame}{What this lab teaches you}
\vspace{0.8em}
\begin{itemize}
  \item The \textbf{ATT is fixed} ($\approx +\$1{,}794$). \;The controls only change the bias.
  \item \textbf{Additive $X$ barely helps:} \$3{,}522 vs naive \$3{,}621. \;You need a design-aware estimator.
  \item OR, IPW, DR-IPW, DR-improved offer four different uses of the same $X$ --- all land within \$240 of the benchmark.
  \item On this dataset, \textbf{OR (regression-adjusted) lands closest} at \$24 below.
  \item No method is magic. Each rests on an assumption you can name.
\end{itemize}
\end{frame}

%% =====================================================================
%%                                CLOSING
%% =====================================================================
\begin{frame}[plain]
\begin{tikzpicture}[remember picture,overlay]
\fill[cream] (current page.south west) rectangle (current page.north east);
\node[anchor=center, text width=0.84\paperwidth, align=center, text=charcoal] at (current page.center) {%
  \fontsize{20pt}{30pt}\selectfont
  The ATT lives on the \textbf{treated.}\\[0.6em]
  The controls only tell you \\[0.2em]
  whether your method \\[0.2em]
  can {\color{accent}\bfseries see it.}
};
\draw[accent, line width=1pt] ([yshift=1.8cm,xshift=0.28\paperwidth]current page.south west) -- ([yshift=1.8cm,xshift=-0.28\paperwidth]current page.south east);
\end{tikzpicture}
\end{frame}

\end{document}
